\begin{table}[H]
\caption{Tipos de retroalimentación y su efecto en los estudios que responden a RQ4.\label{tab:rq4-feedback}}
%\isPreprints{\centering}{%
\begin{adjustwidth}{-\extralength}{0cm}
%}
%\isPreprints{\begin{tabularx}{\textwidth}{CCCC}}{%
\begin{tabularx}{\fulllength}{p{0.25\fulllength} p{0.31\fulllength} X p{0.06\fulllength}}
%}
\toprule
\textbf{Tipo de retroalimentación} & \textbf{Destino / efecto} & \textbf{Estudios designados} & \textbf{Total} \\
\midrule

Resultados de validación o ejecución & Corrección o regeneración de ontologías, modelos y artefactos generados. & \cite{Sabaliauskaite2026,MarianPasca2025,Gstur2024} & 3 \\

Evaluación de expertos o desarrolladores & Revisión de requisitos, casos reutilizados y decisiones de refactoring. & \cite{Jung2023,Bacchiega2024} & 2 \\

Telemetría, rendimiento y cumplimiento de SLA & Optimización de configuración, recursos, políticas y decisiones posteriores. & \cite{William2025,Yamsani2026} & 2 \\

Guardrails y autoevaluación del modelo & Reparación de candidatos arquitectónicos o decisión sobre el uso de conocimiento recuperado. & \cite{Salem2026,Guan2024} & 2 \\

Relevancia de consulta y estado de la base de conocimiento & Restricción y actualización dinámica del espacio de conocimiento utilizado en la recuperación. & \cite{Jin2025} & 1 \\

\midrule
\textbf{Total} & & & \textbf{10} \\
\bottomrule
\end{tabularx}
%\isPreprints{}{%
\end{adjustwidth}
%}
\end{table}